\unitopen{30}
\begin{candos}
\item czytać i rozumieć teksty narracyjne w Präteritum (artykuły, biografie, opowiadania);
\item znać Präteritum najczęstszych czasowników regularnych i nieregularnych (\de{ging, kam, sah, dachte});
\item opowiedzieć krótką historię z \de{Plusquamperfekt} (\de{hatte} / \de{war} + PII);
\item łączyć zdarzenia przysłówkami (\de{zunächst, dann, plötzlich, schließlich}).
\end{candos}

\section{Słownictwo}
\vocabtable{30}

\section{Präteritum: forma pisana i narracja}
Regularne: rdzeń + \de{-te}: \de{machen $\rightarrow$ machte}, \de{arbeiten $\rightarrow$ arbeitete}. Nieregularne mają zmienioną samogłoskę rdzenia;
koniecznie ucz się ich w trzech formach (bezokolicznik, Präteritum, Partizip II).

\begin{center}\small
\begin{tabular}{@{}llll@{}}
\toprule
\textsf{\textbf{Bezokolicznik}} & \textsf{\textbf{ich/er (Prät.)}} & \textsf{\textbf{wir/sie (Prät.)}} & \textsf{\textbf{Partizip II}}\\
\midrule
\de{gehen} & \de{ging} & \de{gingen} & \de{gegangen}\\
\de{kommen} & \de{kam} & \de{kamen} & \de{gekommen}\\
\de{sehen} & \de{sah} & \de{sahen} & \de{gesehen}\\
\de{denken} & \de{dachte} & \de{dachten} & \de{gedacht}\\
\de{bringen} & \de{brachte} & \de{brachten} & \de{gebracht}\\
\de{finden} & \de{fand} & \de{fanden} & \de{gefunden}\\
\de{schreiben} & \de{schrieb} & \de{schrieben} & \de{geschrieben}\\
\de{sprechen} & \de{sprach} & \de{sprachen} & \de{gesprochen}\\
\bottomrule
\end{tabular}
\end{center}
\textbf{Plusquamperfekt} (zdarzenie wcześniejsze niż inne w przeszłości): \de{hatte}/\de{war} + Partizip II:
\de{Nachdem ich die Prüfung bestanden \textbf{hatte}, feierte ich.} \de{Er \textbf{war} schon \textbf{gegangen}, als ich ankam.}

\begin{zasada}[Perfekt w mowie, Präteritum w tekście]
Opowiadając w rozmowie, użyj Perfektu. W piśmie narracyjnym (prasa, biografie, powieści) dominuje Präteritum: \de{Er wurde 1980 in Bern
geboren, studierte Informatik und zog 2005 nach Zürich.} Musisz je \emph{rozumieć}, a używać aktywnie głównie dla
\de{sein}, \de{haben} i modalnych.
\end{zasada}

\begin{falle}
Polski nie ma wyraźnego Präteritum, więc często dodajesz \de{ge-}: \de{Ich gegang} jest błędem; Präteritum \de{ging} bez
\de{ge-} i bez końcówki dla \de{ich} i \de{er}. Drugi: nie mieszaj czasów w jednym akapicie bez powodu. Trzeci: \de{als} (nie \de{wenn})
dla jednorazowego zdarzenia w narracji: \de{Als ich ankam, regnete es.}
\end{falle}

\begin{szwajcaria}
Wiadomości SRF, artykuły NZZ i biografie w CV (\de{Er wurde \dots{} geboren}) używają Präteritum. W Szwajcarii ludzie
rzadziej mówią Präteritum niż w Niemczech, więc usłyszysz Perfekt. Do ćwiczenia czytania wybieraj krótkie biografie i wiadomości
lokalne.
\end{szwajcaria}

\begin{praca}
Postmortem: \de{Um 14 Uhr fiel der Server aus. Zunächst bemerkte niemand das Problem. Dann schlug das Monitoring Alarm.
Schließlich fanden wir die Ursache: ein Zertifikat war abgelaufen.} Ten rodzaj tekstu piszesz w Präteritum.
\end{praca}

\section{Dialogi}
\textbf{Dialog 1 --- opowiadanie o wydarzeniu}
\begin{dialog}
\turn{A}{Wie war dein erster Arbeitstag in Zürich?}{Jak minął twój pierwszy dzień pracy w Zurychu?}
\turn{B}{Zuerst kam ich zu spät, weil der Zug Verspätung hatte. Dann sprach ich mit dem Team.}{Najpierw spóźniłem się, bo pociąg miał opóźnienie. Potem rozmawiałem z zespołem.}
\turn{A}{Und was passierte danach?}{I co stało się potem?}
\turn{B}{Plötzlich fiel das Netzwerk aus. Wir fanden das Problem erst am Abend.}{Nagle padła sieć. Znaleźliśmy problem dopiero wieczorem.}
\end{dialog}

\section{Ćwiczenia}
\begin{cwiczenia}
\exitem{Präteritum: \de{gehen} (er).}{\de{ging}}
\exitem{Präteritum: \de{kommen} (ich).}{\de{kam}}
\exitem{Präteritum: \de{sehen} (wir).}{\de{sahen}}
\exitem{Präteritum: \de{arbeiten} (sie, ona).}{\de{arbeitete}}
\exitem{Präteritum: \de{finden} (ich).}{\de{fand}}
\exitem{Plusquamperfekt: \de{Er ist schon gegangen. (wcześniej niż moje przyjście)}}{\de{Er war schon gegangen, als ich ankam.}}
\exitem{Połącz: \de{Nachdem ich die Prüfung \dots{} (bestehen, Plusquamperfekt), feierte ich.}}{\de{Nachdem ich die Prüfung bestanden hatte, feierte ich.}}
\exitem{Popraw: \de{Ich gegang nach Hause.}}{\de{Ich ging nach Hause.}}
\exitem{Przekształć z Perfekt na Präteritum (pisemne): \de{Er hat 2005 in Zürich gearbeitet.}}{\de{Er arbeitete 2005 in Zürich.}}
\exitem{\de{als} czy \de{wenn}: \dots{} \de{er ankam, regnete es.}}{\de{Als}}
\exitem{Przetłumacz: „Nagle serwer przestał działać.”}{\de{Plötzlich fiel der Server aus.}}
\exitem{Przetłumacz: „Studiował informatykę i przeprowadził się do Zurychu.”}{\de{Er studierte Informatik und zog nach Zürich.}}
\end{cwiczenia}

\weekend{Napisz w Präteritum krótki życiorys fikcyjnej osoby (osiem zdań) i opowiedz go w Perfekcie na głos.
Zaznacz różnice. Przeczytaj krótką biografię z \de{SRF} lub \de{Deutsche Welle} i wypisz wszystkie czasowniki w Präteritum.}
